%!TEX root = ../template.tex
%%%%%%%%%%%%%%%%%%%%%%%%%%%%%%%%%%%%%%%%%%%%%%%%%%%%%%%%%%%%%%%%%%%%
%% abstract-en.tex
%% NOVA thesis document file
%%
%% Abstract in English
%%%%%%%%%%%%%%%%%%%%%%%%%%%%%%%%%%%%%%%%%%%%%%%%%%%%%%%%%%%%%%%%%%%%

\typeout{NT FILE abstract-en.tex}%

Conflict-Free Replicated Data Types (CRDTs) are a key abstraction for modern distributed systems, enabling concurrent
updates across multiple replicas while guaranteeing Strong Eventual Consistency without centralized coordination.
Despite these guarantees, designing CRDTs correctly remains challenging, and formal verification plays an important
role in preventing errors that compromise data integrity. Automated verification tools based on SMT solvers, such as
VeriFx, enable rapid prototyping and the validation of algebraic convergence properties, but face significant limitations
when reasoning about recursive definitions, inductive properties, and global state invariants. Deductive platforms such
as Why3 overcome these limitations through explicit contracts and invariants, at the cost of a much higher specification
effort from the programmer.

This dissertation addresses this trade-off with an integrated framework for the modular, formal development of CRDTs.
Instead of translating directly between VeriFx and Why3, the framework introduces a dedicated language for specifying
state-based and operation-based CRDTs. Once validated through lexical, syntactic, and type analysis over a common abstract
syntax tree, a single specification is automatically compiled into both VeriFx, targeting automated proof of convergence
properties via the Z3 solver, and WhyML, targeting deductive verification of semantic correctness and global invariants.
This design hides the syntactic and modelling complexity of both target ecosystems from the programmer, while ensuring
that the two verification efforts remain consistent with one another, since both are derived from the same specification.
The framework was evaluated on ten CRDTs, including counters, sets, a generic map, and the Replicated Growable Array,
consistently reducing the manual specification effort required while producing idiomatic, tool-specific artifacts for
both targets. In particular, the Replicated Growable Array case study exposes a scenario in which the deductive Why3 proof
succeeds while the fully automated VeriFx/Z3 proof fails, illustrating the benefit of pairing automated and deductive
verification rather than relying on either in isolation.

% Abstract keywords
\keywords{
  CRDTs \and
  Formal Verification \and
  Why3 \and
  VeriFx
}
